\documentclass{article}
\usepackage{amsmath, amssymb, amsthm}
\usepackage{hyperref}
\usepackage{geometry}
\geometry{margin=1in}

\title{Erd\H{o}s Problem \#124}
\date{Last updated: 2025-08-31}
\author{Source: erdosproblems.com}

\begin{document}
\maketitle

\noindent\textbf{Status:} OPEN \quad \textbf{No prize}

\noindent\textbf{Tags:} number theory, base representations

\medskip
\noindent\textbf{Problem Statement:}

For any $d\geq 1$ and $k\geq 0$ let $P(d,k)$ be the set of integers which are the sum of distinct powers $d^i$ with $i\geq k$. Let $3\leq d_1<d_2<\cdots <d_r$ be integers such that\[\sum_{1\leq i\leq r}\frac{1}{d_r-1}\geq 1.\]Can all sufficiently large integers be written as a sum of the shape $\sum_i c_ia_i$ where $c_i\in \{0,1\}$ and $a_i\in P(d_i,0)$? 

If we further have $\mathrm{gcd}(d_1,\ldots,d_r)=1$ then, for any $k\geq 1$, can all sufficiently large integers be written as a sum of the shape $\sum_i c_ia_i$ where $c_i\in \{0,1\}$ and $a_i\in P(d_i,k)$?

\bigskip
\noindent\textbf{Remarks:}

The second question was conjectured by Burr, Erd\H{o}s, Graham, and Li \cite{BEGL96}, who proved it for $\{3,4,7\}$.

The first question was asked separately by Erd\H{o}s in \cite{Er97} and \cite{Er97e} (although there is some ambiguity over whether he intended $P(d,0)$ or $P(d,1)$ - certainly he mentions no gcd condition). A simple positive proof of the first question was provided (and formalised in Lean) by Aristotle thanks to Alexeev; see the comments for details.

In \cite{BEGL96} they record that Pomerance observed that the condition $\sum 1/(d_i-1)\geq 1$ is necessary (for both questions), but give no details. Tao has sketched an explanation in the comments. It is trivial that $\mathrm{gcd}(d_1,\ldots,d_r)=1$ is a necessary condition in the second question. 

Melfi \cite{Me04} gives a construction, for any $\epsilon>0$, of an infinite set of $d_i$ for which every sufficiently large integer can be written as a finite sum of the shape $\sum_i c_ia_i$ where $c_i\in \{0,1\}$ and $a_i\in P(d_i,0)$ and yet $\sum_{i}\frac{1}{d_i-1}<\epsilon$.

See also [125].

\bigskip
\noindent\textbf{References:}

[BEGL96] Burr, S. A. and Erd\H{o}s, P. and Graham, R. L. and Li, W. Wen-Ching, Complete sequences of sets of integer powers. Acta Arith. (1996), 133-138.
 
 [Er97] Erd\H{o}s, Paul, Problems in number theory. New Zealand J. Math. (1997), 155-160.
 
 [Er97e] Erd\H{o}s, Paul, Some of my favourite unsolved problems. Math. Japon. (1997), 527-537.
 
 [Me04] Melfi, Giuseppe, On certain positive integer sequences. Riv. Mat. Univ. Parma (7) (2004), 253--260.

\bigskip
\noindent\small{Source: \url{https://www.erdosproblems.com/124}}
\end{document}
